\documentclass[11pt,a4paper]{article}
\usepackage[margin=2.2cm]{geometry}
\usepackage{booktabs}
\usepackage{array}
\usepackage[T1]{fontenc}
\usepackage{lmodern}
\usepackage{microtype}
\usepackage{xcolor}
\usepackage[colorlinks=true,linkcolor=black,urlcolor=blue!55!black]{hyperref}
\usepackage{titlesec}

\titleformat{\section}{\normalfont\large\bfseries}{\thesection.}{0.6em}{}
\titlespacing*{\section}{0pt}{1.4ex plus .3ex}{0.7ex}
\setlength{\parskip}{0.55em}
\setlength{\parindent}{0pt}
\newcommand{\tc}[1]{\texttt{\small #1}}

\title{\vspace{-1.6cm}\textbf{Can AI Detect XSS in Web Traffic?}\\[0.2em]
\large A guide to what we built, how it works, and what it is for}
\author{\normalsize Chandra Prakash Bathula}
\date{\normalsize \today}

\begin{document}
\maketitle
\vspace{-0.8cm}

\section{What this project actually is}

This is a measurement experiment, not a product. The question we set out to answer is
whether a machine learning detector placed in network traffic can still recognise
cross-site scripting attacks after someone rewrites the payload into a different
representation, and what that detector costs in response time when it runs on every
request.

The important thing to understand up front is that we were never trying to build a good
XSS detector. We were trying to measure one honestly. A detector that scores well is not
the goal. A comparison that is fair, repeatable, and clear about its own limits is the
goal. Everything described below exists to serve that.

All of it runs on the local machine. The only payload used anywhere is a harmless flag,
\tc{window.\_\_xss\_fired=true}, which sets a variable and does nothing else. No real
exploit code is involved and nothing reaches the network.

\section{What is running, and on which port}

There are four ports in play. Three of them are servers you start yourself, and the fourth
only appears while the timing benchmark is running.

\begin{center}
\begin{tabular}{@{}l l p{8.4cm}@{}}
\toprule
\textbf{Port} & \textbf{Process} & \textbf{What it does} \\
\midrule
5000 & \tc{app/vuln\_app.py} &
The deliberately vulnerable web application. This is the thing being attacked. It has five
endpoints, and each one puts the same user input into a different place in the HTML, so the
browser treats it differently. It has no user interface on purpose. \\
\addlinespace
8080 & \tc{proxy/proxy.py} &
The reverse proxy. Every request passes through here before reaching the app, and this is
where the detector lives. It has three modes: \tc{off} forwards without scoring,
\tc{observe} scores and logs but always forwards, and \tc{block} returns 403 when the score
crosses the threshold. \\
\addlinespace
5001 & \tc{teach/playground.py} &
The only part with a real interface. It has two pages. The playground lets you type an
input and watch the detector score it and the browser run it side by side. The results page
shows every measurement table and the figure. \\
\addlinespace
8090 & \tc{bench/loadtest.py} &
A temporary proxy that the benchmark starts and stops by itself, once per run. It uses a
separate port so the benchmark never disturbs whatever is on 8080. \\
\bottomrule
\end{tabular}
\end{center}

The five endpoints on port 5000 are the heart of the experiment. \tc{/reflect-text} escapes
the input before printing it, so it is safe. \tc{/reflect-html} drops it straight into the
page body. \tc{/reflect-attr} puts it inside an HTML attribute. \tc{/reflect-js} puts it
inside a JavaScript string. \tc{/reflect-decode} decodes the input a second time before
printing it, which lets us study what happens when the proxy and the application disagree
about how many times to decode something.

The same input sent to \tc{/reflect-text} and to \tc{/reflect-html} produces an identical
HTTP request. The detector therefore gives it an identical score. One of them is harmless
and the other one runs code. That gap is the entire subject of the project.

\section{What we built and why}

The vulnerable app came first, and it is kept as plain as possible. If we gave it a nice
interface we would be measuring the interface instead of the detector. It is a test fixture
and it should stay boring.

The proxy came second. It matters that the detector sits here, because this is the only
place a network defence can realistically live. It can read the request line, the headers,
and the query string. It cannot see what the application will do with the value, and it
cannot see whether the browser ends up running anything. We log both the raw query string
and the decoded value on every request so that failures can be explained afterwards.

The dataset came third, and it took the longest. We collected 1000 labelled inputs, half
attack-like and half harmless, from four sources: a Kaggle XSS dataset, the PayloadBox
payload list as mirrored in SecLists, the benign rows of a large Mendeley dataset, and 50
tricky harmless strings we wrote by hand. The hand-written ones matter because they contain
things like \tc{<b>bold</b>} and \tc{if (a < b \&\& c > d)}, which look like attacks to a
naive detector but are perfectly ordinary text.

The detector itself is deliberately simple. It breaks text into character sequences of
length two to five, weights them, and runs logistic regression. The baseline model ends up
with 13{,}617 of these features. The decision threshold, 0.617, was chosen on a validation
set and then frozen. We never changed the model or the features afterwards, because the
experiment compares data, not algorithms.

\section{How we did it, step by step}

We built the app and the proxy first and proved that a request through the proxy produced
byte-for-byte the same page as a request straight to the app. Without that, nothing
measured later would mean anything.

Then we cleaned and split the data. Public payload lists are full of near-copies that
differ by one digit, so we grouped inputs by a skeleton key that replaces numbers and
quoted strings with placeholders. Every variant we later generated inherits its parent's
group and its parent's split. The code asserts that no group ever appears in two splits. If
that assertion ever fires, the results are worthless, so it runs every time.

Next we trained the baseline and recorded how it does on ordinary inputs. It does very
well, with an F1 of 0.995. That number is supposed to be high. It is the number the rest of
the experiment knocks down.

Then we wrote the three transformation families. The first re-encodes the payload for
transport using percent encoding. The second rewrites characters as HTML entities. The
third rewrites the JavaScript using bracket notation, string concatenation, and hex
escapes. This produced 5012 variants in total, including 1500 harmless ones used only to
check the false positive rate.

Every attack variant was then opened in a real Chromium browser at its intended endpoint,
and we recorded whether the marker actually fired. This gave us two separate labels that we
are careful never to mix. One says the input looks like an attack. The other says the
browser really ran it. A great many inputs satisfy the first and fail the second.

We also found that 56 percent of the payloads only fire when a user clicks or hovers. We
record that fact rather than firing the events ourselves, because forcing a click would
quietly change the meaning of the measurement.

Finally we retrained the same model with extra training data, holding the JavaScript family
out of both training and validation, and measured the cost of running the detector inline
under two traffic loads, three times each, with the configuration order alternated between
rounds.

\section{What we found}

On ordinary inputs the detector catches 95 percent of attacks. On the same attacks after
transformation it catches 60 percent. The confidence intervals do not overlap, so the drop
is real and not an artefact of which rows landed in the test set.

The sharper result comes from comparing the two labels. Of the variants that genuinely
executed in the browser, the detector flagged almost none, 0 percent for the entity family
and 5 percent for the JavaScript family. Of the attack-shaped inputs that did nothing at
all, it flagged between 79 and 98 percent. In other words it reacts to how text looks
rather than to what it does.

Augmenting the training data appeared to fix everything, with a perfect score in every
condition. We did not believe it, so we renamed the marker identifiers in the test inputs
while leaving the structure and the obfuscation untouched. The score on the held-out family
fell from 1.000 to 0.726, which is level with the baseline's 0.709. The improvement was
mostly the model learning our marker's name. Its false positive rate also rose from 0.010
to 0.023. This is a negative result, and the brief says plainly that a negative result is
valuable when it explains itself.

The cost of running the detector is small. It adds about half a millisecond to the median
response time. Throughput settles at roughly 711 requests per second with the detector on
and also with it off, which means the limit is the application behind it rather than the
model.

\section{What all of this is useful for}

The direct use is that the project answers its own research question with evidence instead
of opinion. It shows why the detector misses transformed inputs, which is that it learned
surface strings and not behaviour. The strongest features it learned are literally the
letters in the word \tc{alert} followed by an opening bracket.

The wider use is the conclusion it supports. A detector reading HTTP requests cannot tell
whether the value it is looking at will be escaped safely or dropped into a script tag,
because that decision happens later and somewhere else. Our paired tests make this concrete.
Eight inputs, each sent to a safe endpoint and an unsafe one, produced identical scores and
opposite outcomes every time. The safe endpoint was safe because it escaped its output, not
because anything was detected.

That is the practical lesson worth carrying out of this work. Detection of this kind is
worth having as a signal and as a source of logs, and it is cheap enough to run inline. It
is not a replacement for encoding output correctly in the application. Anyone who treats a
request-content classifier as the defence has misunderstood where the vulnerability lives.

The playground on port 5001 exists to make that argument in about thirty seconds. Load the
entity-encoded example, watch the detector return a verdict of ALLOW, and watch the browser
run the marker anyway in the panel directly underneath.

\end{document}
